% !TeX root = ../main.tex

\chapter{Conclusion}\label{chapter:conclusion}

This thesis demonstrated that an Incremental Nonlinear Dynamic Inversion (INDI) controller can maintain stable and accurate flight on a quadrotor whose aerodynamic surfaces generate lift on the same order as its weight—without using any explicit aerodynamic model.

A fully 3D-printed X-wing quadrotor was designed, modeled, and flight-tested.
The static thrust mapping followed a quadratic law with $R^2=0.997$, confirming consistent actuator behavior.
During motion-capture experiments, the controller achieved RMS position errors between 0.03\,m and 0.16\,m up to 8\,m\,s$^{-1}$ forward speed, demonstrating precise trajectory tracking under strong, unmodeled aerodynamic disturbances.

Outdoor experiments with direct electrical-power logging confirm that aerodynamic lift yields tangible energy savings. Power consumption decreased by approximately 33\% at 10\,m\,s$^{-1}$ relative to hover, validating that passive lift directly reduces rotor workload.
The progression from the initial $45^{\circ}$ NACA~0015 design to the refined $20^{\circ}$ AG25 variant demonstrates a direct link between aerodynamic optimization and measurable electrical-power reduction.
The improved wing geometry validated under outdoor conditions achieved this 33\% power reduction at 10\,m\,s$^{-1}$ compared with hover.
These findings close the loop between aerodynamic theory, thrust-based proxy estimates, and real-flight measurements, providing the first quantitative evidence of efficiency improvement in an INDI-controlled aerodynamic multirotor.

The results confirm that disturbance-rejection control can absorb quasi-steady aerodynamic effects without identification or model compensation.
This finding simplifies control design for hybrid multirotors: aerodynamic augmentation can be pursued purely at the hardware level, leaving the control law unchanged.

The main limitations are the absence of direct power measurements and quantitative comparison to a geometric-only baseline.
Future work should instrument the system with current sensing to quantify energy savings, and perform controlled A/B tests to isolate the INDI improvement.
Extending the framework to dynamic flight envelopes and real-time aerodynamic estimation would further test its generality.
Additional future work includes:
\begin{itemize}
\item Extend tests to longer endurance missions to quantify range increase and include varying payload configurations.
\end{itemize}

Overall, the study establishes a reproducible experimental benchmark for aerodynamic surface-enhanced multirotors and provides evidence that model-free incremental control remains robust even when aerodynamic lift dominates vehicle dynamics.

